\chapter{نکاتِ پیشرفته و بهترین روش‌ها}%
\label{ch:advanced}

به فصلِ پایانیِ بخشِ اصلیِ کتاب رسیدیم. تا اینجا ابزارهای زبان را آموختید؛ اکنون
وقتِ آن است که یاد بگیریم چگونه از آن‌ها \emph{خوب} استفاده کنیم: کدی که سریع،
ایمن، خوانا و قابلِ‌اعتماد باشد. این فصل، صندوقچهٔ ابزارِ یک برنامه‌نویسِ
حرفه‌ای است.

\section{پاس دادن به ارجاع}

در فصلِ ۶ دیدیم که پارامترها معمولاً یک \emph{رونوشت} از مقدار می‌گیرند؛ پس
تغییرِ آن‌ها درونِ تابع، اصلِ متغیرِ بیرون را عوض نمی‌کند. اما گاهی می‌خواهیم
تابع \emph{خودِ} متغیر را تغییر دهد. برای این از \textbf{ارجاع} (با نشانهٔ
\code{\&}) استفاده می‌کنیم. این مثالِ واقعی را ببینید:

\begin{salam}
تابع افزایش(ایکس &: صحیح):
    ایکس = ایکس + 1
تمام

تابع جابجایی(الف &: صحیح, ب &: صحیح):
    متغیر موقت : صحیح = الف
    الف = ب
    ب = موقت
تمام

تابع اصلی:
    متغیر ن : صحیح = 10
    افزایش(ن)
    افزایش(ن)
    چاپ "ن =", ن
    متغیر پ : صحیح = 1
    متغیر ک : صحیح = 2
    جابجایی(پ, ک)
    چاپ "پ =", پ, "ک =", ک
تمام
\end{salam}

پارامترِ \code{ایکس \&صحیح} یک \emph{ارجاع} به متغیرِ اصلی است، نه رونوشتِ
آن. پس وقتی \code{افزایش(ن)} را دو بار صدا می‌زنیم، \code{ن} واقعاً به ۱۲
می‌رسد. به همین شکل، \code{جابجایی} می‌تواند جای دو متغیرِ بیرونی را عوض کند ---
کاری که بدونِ ارجاع ناممکن بود.

\begin{note}
درونِ بدنهٔ تابع، با ارجاع درست مانندِ یک متغیرِ معمولی کار می‌کنید ---
\code{ایکس = ایکس + 1}. نیازی به نشانهٔ ویژه نیست. تنها در \emph{امضای} تابع
است که \code{\&} می‌گوید این پارامتر یک ارجاع است.
\end{note}

\section{قالب‌بندیِ خودکارِ کد}

کدِ مرتب و یکدست، خواندن و نگه‌داری را آسان می‌کند. سلام یک \textbf{قالب‌بند}
(formatter) داخلی دارد که کد را به‌طورِ خودکار و طبقِ سبکِ استانداردِ زبان مرتب
می‌کند:

\begin{salamen}[language={}]
salam format app.salam     # مرتب کردنِ یک فایل
salam format               # مرتب کردنِ همهٔ فایل های پروژه
salam format --check       # فقط گزارش بده کدام فایل ها نامرتب اند
\end{salamen}

\begin{tip}
بحث بر سرِ «چند فاصله تورفتگی» یا «پرانتز کجا باشد» اتلافِ وقت است. بگذارید
\code{salam format} این تصمیم‌ها را بگیرد. تیم‌های حرفه‌ای معمولاً قالب‌بند را
طوری تنظیم می‌کنند که پیش از هر ثبت (commit) خودکار اجرا شود؛ این‌گونه کلِ کدِ
پروژه همیشه یکدست می‌ماند.
\end{tip}

\section{آزمون‌نویسی}

چگونه مطمئن شویم که برنامه‌مان درست کار می‌کند و پس از هر تغییر، همچنان درست
می‌ماند؟ پاسخ، \textbf{آزمون} است. ماژولِ \code{testing} ابزارهای آزمون را
فراهم می‌کند:

\begin{salam}
واردکردن "testing"
تابع جمع(الف : صحیح, ب : صحیح) : صحیح: بازگشت الف + ب تمام

تابع اصلی:
    testing.AssertEqInt(جمع(2, 3), 5)
    testing.AssertEqInt(جمع(-1, 1), 0)
    testing.AssertEqInt(جمع(0, 0), 0)
    testing.Summary()
تمام
\end{salam}

تابعِ \code{AssertEqInt(واقعی, مورِدِانتظار)} بررسی می‌کند که نتیجهٔ تابع همان
چیزی است که انتظار داشتیم. اگر همه درست باشند، آزمون‌ها سبز می‌شوند؛ اگر یکی
شکست بخورد، دقیقاً می‌گوید کجا. \code{Summary()} در پایان خلاصه‌ای از نتیجه
نشان می‌دهد.

\begin{warn}
عادتِ مهلک: «بعداً آزمون می‌نویسم.» در عمل، «بعداً» هرگز نمی‌آید. آزمون‌ها را
\emph{همراهِ} کد بنویسید، نه پس از آن. هر تابعِ مهم باید دستِ‌کم چند آزمون
داشته باشد به‌ویژه برای حالت‌های مرزی (صفر، عددِ منفی، رشتهٔ خالی) که
بیشترین خطاها در آن‌ها رخ می‌دهد.
\end{warn}

\section{اشکال‌زدایی}

وقتی برنامه آن‌طور که می‌خواهید کار نمی‌کند، به ابزارِ \textbf{اشکال‌زدا}
(debugger) نیاز دارید که اجازه دهد گام‌به‌گام برنامه را پیش ببرید و مقدارِ
متغیرها را ببینید. سلام با گزینهٔ \code{-g} اطلاعاتِ اشکال‌زدایی تولید می‌کند:

\begin{salamen}[language={}]
salam debug app.salam      # ساخت با اطلاعاتِ اشکال زدایی و اجرای اشکال زدا
\end{salamen}

افزون بر این، برای یافتنِ نشت‌های حافظه، فرمانِ \code{memcheck} را دارید:

\begin{salamen}[language={}]
salam memcheck app.salam   # اجرای برنامه زیرِ ابزارِ وارسیِ حافظه
\end{salamen}

\begin{tip}
ساده‌ترین ابزارِ اشکال‌زدایی همان \code{چاپ} است! وقتی گیر کردید، چند
\code{چاپ} در نقاطِ کلیدی بگذارید تا ببینید مقدارها چه هستند و برنامه به کجا
می‌رسد. این روشِ ساده، اغلب سریع‌ترین راهِ یافتنِ خطاست. اما برای مشکلاتِ
پیچیده، اشکال‌زدای واقعی ارزشش را دارد.
\end{tip}

\section{کارایی و حافظه}

سلام به سرعتِ \lr{C} نزدیک است، چون برنامهٔ شما در نهایت به کدِ ماشین تبدیل
می‌شود. اما چند نکته به برنامه‌های سریع‌تر و امن‌تر می‌انجامد:

\begin{itemize}
  \item \textbf{نوعِ درست را برگزینید:} برای شمارنده‌های کوچک \code{صحیح}
        کافی است؛ \code{صحیح۶۴} را فقط وقتی لازم است به‌کار ببرید.
  \item \textbf{منابع را آزاد کنید:} بردارها و نگاشت‌هایی که عمرشان تمام شده را
        با \code{.آزاد()}/\code{.free()} آزاد کنید. (به یاد دارید که
        \code{تعلیق} بهترین جا برای این کار است؟)
  \item \textbf{از بازگشتِ بیش‌از‌اندازه عمیق بپرهیزید:} هر فراخوانیِ بازگشتی
        فضای پشته می‌گیرد. برای عمقِ زیاد، گاهی حلقه بهتر از بازگشت است.
  \item \textbf{در حلقه‌های داغ، کارِ تکراری را بیرون بکشید:} اگر چیزی در هر
        تکرارِ حلقه ثابت است، آن را \emph{پیش} از حلقه حساب کنید.
\end{itemize}

\begin{deep}[نگاهی به درون: امنیتِ حافظه]
بسیاری از خطاهای جدی و آسیب‌پذیری‌های امنیتی از سوءِاستفادهٔ حافظه می‌آیند:
دسترسی بیرون از مرزِ آرایه، استفاده از حافظه‌ای که آزاد شده، یا فراموشیِ
آزادسازی (نشتِ حافظه). سلام چند سپر در برابرِ این‌ها دارد: در حالتِ توسعه
(پیش‌فرض) دسترسیِ آرایه را وارسی می‌کند، \code{memcheck} نشت‌ها را می‌یابد، و
\code{تعلیق} آزادسازی را تضمین می‌کند. در حینِ توسعه از این ابزارها بهره بگیرید
تا خطاها را \emph{زود} پیش از رسیدن به دستِ کاربر بگیرید؛ برای ساختِ نهایی،
گزینهٔ \code{--release} همین وارسیِ آرایه را برای سرعتِ بیشتر حذف می‌کند.
\end{deep}

\section{اصول‌ِ کدِ خوب}

در پایان، چند اصلِ کلی که برنامه‌نویسانِ باتجربه به آن‌ها پایبندند:

\begin{enumerate}
  \item \textbf{نام‌های گویا:} \code{مجموعِ‌نمرات} بهتر از \code{م} است. کد را
        برای \emph{انسان} بنویسید.
  \item \textbf{توابعِ کوچک:} هر تابع یک کارِ روشن انجام دهد. اگر تابعی خیلی
        بزرگ شد، احتمالاً باید آن را بشکنید.
  \item \textbf{از تکرار بپرهیزید:} اگر کدی را دو بار نوشتید، آن را به یک تابع
        تبدیل کنید.
  \item \textbf{تغییرناپذیری را پیش‌فرض بگیرید:} متغیرها را تا وقتی واقعاً نیاز
        به تغییر نیست، تغییرناپذیر بگذارید.
  \item \textbf{خطاها را آگاهانه مدیریت کنید:} هیچ حالتِ خطایی را نادیده
        نگیرید.
  \item \textbf{آزمون بنویسید:} کدِ بدونِ آزمون، کدِ مشکوک است.
  \item \textbf{توضیحِ «چرا»، نه «چه»:} توضیحات باید بگویند \emph{چرا} کاری
        انجام می‌شود، نه \emph{چه} کاری چون «چه» را خودِ کد می‌گوید.
\end{enumerate}

\begin{tip}
بهترین کد، کدی نیست که هوشمندانه‌ترین باشد، بلکه کدی است که \emph{ساده‌ترین} و
\emph{خواناترین} باشد. کسی که شش ماه بعد کدِ شما را می‌خواند که اغلب خودِ
شمایید! باید بتواند آن را به‌آسانی بفهمد. سادگی، بالاترین هنرِ
برنامه‌نویسی است.
\end{tip}

\section{راهِ پیشِ رو}

این کتاب پایان یافت، اما سفرِ شما تازه آغاز شده است. برنامه‌نویسی مهارتی است که
تنها با \emph{تمرین} به‌دست می‌آید. پیشنهادهایی برای ادامه:

\begin{itemize}
  \item پروژه‌های کوچکِ خود را بسازید: یک ماشین‌حساب، یک بازیِ حدسِ عدد، یک
        دفترچهٔ یادداشت.
  \item کدِ مثال‌های پوشهٔ \code{examples} سلام را بخوانید و تغییر دهید.
  \item هر مفهومی را که یاد گرفتید، با نوشتنِ یک برنامهٔ کوچک بیازمایید.
  \item از خطاها نترسید؛ هر خطا یک فرصتِ یادگیری است.
\end{itemize}

\section{خلاصهٔ فصل}

\begin{itemize}
  \item با ارجاع (\code{\&}) می‌توان متغیرِ بیرون را از درونِ تابع تغییر داد.
  \item \code{salam format} کد را خودکار مرتب می‌کند.
  \item ماژولِ \code{testing} برای آزمون‌نویسی است.
  \item \code{salam debug} و \code{memcheck} برای اشکال‌زدایی و یافتنِ نشتِ
        حافظه‌اند.
  \item نوعِ درست، آزادسازیِ منابع و پرهیز از تکرار، کلیدهای کارایی‌اند.
  \item کدِ خوب، کدِ ساده و خواناست.
\end{itemize}

\section{تمرین‌ها}

\exercise تابعِ \code{جابجایی} را برای دو رشته بازنویسی کنید (با ارجاع).

\exercise چند آزمون برای تابعِ \code{فاکتوریل}ِ فصلِ ۶ بنویسید، از جمله برای
حالت‌های مرزی.

\exercise یک قطعه‌کدِ نامرتب بنویسید، آن را با \code{salam format} مرتب کنید و
تفاوت را ببینید.

\exercise یکی از برنامه‌های فصل‌های پیش را بردارید و آن را با نام‌های گویاتر و
توابعِ کوچک‌تر بازنویسی کنید.

\exercise یک پروژهٔ کوچکِ کامل طراحی کنید (مثلاً یک دفترِ تلفن با نگاشت) که از
دستِ‌کم پنج مفهومِ این کتاب استفاده کند، و برایش آزمون بنویسید.
